% SHARD-ID: NB-02-REPRO-MAP
% SHARD-TITLE: Reproducibility Map
% SHARD-SUMMARY: Maps each simulation to the runs it produced and each run to the claims and figures that depend on it.
% SHARD-SUMMARY: States the rules that make a number in this lab book regenerable from a recorded command.
% SHARD-KEYWORDS: reproducibility, runs, simulations, figures, artefacts, provenance

\section{Reproducibility Map}
\label{sec:reproducibility-map}

\subsection*{Simulations and their runs}

No simulation has produced a run yet. Once one has, this table maps each simulation to the runs it
produced and each run to the claims that depend on it, per \texttt{labbook/README.md} rule~6.

\subsection*{The contract}

Every number and figure in this lab book must be regenerable. Concretely, a \evC{run-id} tag points
at \texttt{runs/<run-id>/README.md}, and that file records:

\begin{itemize}
  \item the hypothesis the run was meant to test;
  \item the simulation and the git commit SHA it was run at;
  \item the exact command, and the parameters as a committed \texttt{params.toml};
  \item the environment, identified by the \texttt{pixi.lock} hash;
  \item the headline result, its interpretation, and its uncertainty --- convergence level,
        tolerances, statistical error, and the physics the model omits;
  \item the gates that were run.
\end{itemize}

Large outputs are not committed. Each is represented in its run README by a stub giving path, size,
SHA256 and the command that produces it, so its absence is auditable and its regeneration is
mechanical. Small text summaries and the vector figures that appear here are committed.

\subsection*{Run to shard map}

No run has been promoted into a shard yet.

\subsection*{Validation targets}

Every simulation must be checked against something outside itself before its output is used in an
argument: an analytic limit, a published number, a reproduced figure, or a convergence study.
``It ran without errors'' is not a check. Validation targets accepted for each simulation are
listed in that simulation's own README and repeated here once a run has confirmed them. None has
run yet.
